\section*{अध्याय \ref{ch:b3:representations} --- परिमित समूहों के
निरूपण}

\begin{solution}{exo:b3:representations:1}
(क) $\Z/n\Z$ आबेली है: अतः सभी अखंडनीयों की घात $1$ है (\cref{prop:b3:representations:onedim}), अर्थात्
वे समाकारिताएँ $\chi \colon \Z/n\Z \to \C^\times$ हैं, जो $\omega^n = 1$ वाले $\chi(\bar 1)
= \omega$ से निर्धारित होती हैं:
अर्थात् $n$ \omterm{def:b3:representations:character}{वर्ण} $\chi_k(\bar
m) = \eu^{2\iu\pi km/n}$। $n = 4$ के लिए (वर्ग ${}={}$ अवयव $\bar0,\bar1,\bar2,\bar3$):
\[
\begin{array}{c|cccc}
& \bar0 & \bar1 & \bar2 & \bar3\\
\hline
\chi_0 & 1 & 1 & 1 & 1\\
\chi_1 & 1 & \iu & -1 & -\iu\\
\chi_2 & 1 & -1 & 1 & -1\\
\chi_3 & 1 & -\iu & -1 & \iu
\end{array}
\]
(ख) पंक्तियाँ: $\langle\chi_k,\chi_l\rangle = \frac14\sum_m
\eu^{2\iu\pi(l-k)m/4} = \delta_{kl}$ (ज्यामितीय योग); और स्तंभ भी वैसे ही। आव्यूह
$(\eu^{2\iu\pi km/n})_{k,m}$ \emph{विविक्त फूरिये रूपांतर} आव्यूह है --- वही
इकाई-मूल छन्ना जो दूसरे वर्ष के खंड के जनक-फलनों वाले अध्याय में
प्रयुक्त हुआ था; \omterm{thm:b3:representations:orthogonality}{वर्णों की लांबिकता} विविक्त फूरिये रूपांतर के
प्रतिलोमन सूत्र का सामान्यीकरण है।
\end{solution}

\begin{solution}{exo:b3:representations:2}
(क) आधार $(e_x)$ में $\rho(g)$ का आव्यूह कोई क्रमचय आव्यूह है, जिसका संचिह्न
$g\cdot x =
x$ वाले $x$ की संख्या है। तब बर्नसाइड की गणना प्रमेयिका (\cref{exo:b3:groups:5})
से
\[
\langle\mathbf 1, \chi\rangle = \frac1{\abs G}\sum_g
\abs{\operatorname{Fix}(g)} = \#\{\text{कक्षाएँ}\}
\]
--- समतुल्य रूप से, यह तुच्छ \omterm{def:b3:representations:rep}{निरूपण} की बहुलता परिकलित करता है, जिसका
समजातीय स्थान $G$-निश्चर सदिशों का स्थान है, और उसकी विमा कक्षाओं
की संख्या है (प्रति \omterm{def:b3:groups:action}{कक्षा} एक सूचक)।

(ख) चूँकि $\operatorname{Fix}_{X\times X}(g) =
\operatorname{Fix}_X(g)^2$, अतः $X \times X$ पर लगाया गया भाग (क) $\langle\chi,\chi\rangle = \frac1{\abs
G}\sum\abs{\operatorname{Fix}(g)}^2 = \#\{
X\times X \text{ पर कक्षाएँ}\}$ दे देता है
($\chi$ वास्तविक है)। $\chi = \mathbf 1 +
\psi$ लिखने पर: $\langle\mathbf1,\chi\rangle = 1$ (संक्रमणीयता), अतः $\langle\psi,\psi\rangle = \langle\chi,\chi\rangle - 1$।
$X\times X$ पर क्रिया में विकर्ण एक \omterm{def:b3:groups:action}{कक्षा} है; और ठीक एक अन्य \omterm{def:b3:groups:action}{कक्षा} होती
है तभी और केवल तभी जब $G$ भिन्न युग्मों पर संक्रमणीय हो: अतः $\langle\psi,\psi\rangle = 1$
तभी और केवल तभी जब $2$-संक्रमणीय (\cref{cor:b3:representations:multiplicity})।

(ग) $S_n$ $\{1,\dots,n\}$ पर $2$-संक्रमणीय है (किसी भी भिन्न युग्म को कहीं भी
भेजिए): अतः $\psi = \chi_{\mathrm{std}}$ \omterm{def:b3:representations:rep}{अखंडनीय} है।
\end{solution}

\begin{solution}{exo:b3:representations:3}
$S_3$: तीन वर्ग, $\sum n_i^2 = 6 = 1 + 1 + 4$; दोनों घात-$1$ \omterm{def:b3:representations:character}{वर्ण} $\mathbf 1, \varepsilon$ हैं ($[S_3 :
A_3] = 2$); और तीसरी
पंक्ति $(2, a, b)$ $e$-स्तंभ के साथ स्तंभ लांबिकता से निकलती है: $1 - 1 + 2a = 0$ और
$1 + 1 +
2b = 0$: अतः $a = 0$, $b = -1$ --- अर्थात् \cref{ex:b3:representations:s3} की सारणी।

$S_4$ जाँच (पंक्तियाँ): $\langle\chi_{\mathrm{std}},
\varepsilon\chi_{\mathrm{std}}\rangle = \frac1{24}(9 - 6 + 3 + 0
- 6) = 0$; $\langle\chi_2,\chi_2\rangle = \frac1{24}(4 + 0 + 12
+ 8 + 0) = 1$। स्तंभ: $(1\,2)$ के विरुद्ध $e$:
$1 - 1 + 0 + 3 - 3
= 0$; $(1\,2)$ स्वयं के विरुद्ध: $1 + 1 + 0 + 1 + 1 = 4 =
\abs{Z_{S_4}((1\,2))} = 24/6$।

$4$ बिंदुओं पर क्रमचय \omterm{def:b3:representations:character}{वर्ण}: $(4, 2, 0, 1, 0) = \mathbf
1 + \chi_{\mathrm{std}}$ (स्थिर-बिंदु गणनाएँ; शीर्ष पंक्ति
घटाइए)। $\chi_{\mathrm{std}}^2 = (9, 1, 1, 0, 1)$ के लिए:
\[
\langle\mathbf1,\cdot\rangle = \tfrac{9 + 6 + 3 + 0 + 6}{24} =
1,\quad
\langle\varepsilon,\cdot\rangle = 0,\quad
\langle\chi_2,\cdot\rangle = \tfrac{18 + 6}{24} = 1,\quad
\langle\chi_{\mathrm{std}},\cdot\rangle = 1,\quad
\langle\varepsilon\chi_{\mathrm{std}},\cdot\rangle = 1:
\]
$\chi_{\mathrm{std}}^2 = \mathbf 1 + \chi_2 +
\chi_{\mathrm{std}} + \varepsilon\chi_{\mathrm{std}}$
(विमाएँ: $9 = 1 + 2 + 3 + 3$)।
\end{solution}

\begin{solution}{exo:b3:representations:4}
दोनों समूहों के पाँच-पाँच वर्ग हैं और घात प्रतिरूप $(1,1,1,1,2)$ है (आबेलीकरण
$\cong (\Z/2\Z)^2$ से चार घात-$1$, फिर $\sum n_i^2 = 8$)। वर्गों को $e$, $z$ (केंद्रीय
अंतर्वलन: क्रमशः $r^2$ और $-1$), तथा तीनों द्वि-अवयवी वर्गों के क्रम
में रखने पर:
\[
\begin{array}{c|ccccc}
& e & z & C_1 & C_2 & C_3\\
\hline
\chi^{(1)} & 1 & 1 & 1 & 1 & 1\\
\chi^{(2)} & 1 & 1 & 1 & -1 & -1\\
\chi^{(3)} & 1 & 1 & -1 & 1 & -1\\
\chi^{(4)} & 1 & 1 & -1 & -1 & 1\\
\chi^{(5)} & 2 & -2 & 0 & 0 & 0
\end{array}
\]
(अंतिम पंक्ति स्तंभ लांबिकता से)। $D_4$ और $Q_8$ की सारणियाँ समरूप
हैं, जबकि वे तुल्याकारी नहीं हैं (\cref{pb:b3:groups:1})। सारणी जो \emph{पकड़ लेती
है}: $\abs
G$, वर्ग-आकार, \omterm{ex:b3:groups:actions}{केंद्र} ($\{g : \abs{\chi_i(g)} = n_i\
\forall i\}$: दोनों में क्रम $2$), आबेलीकरण,
और प्रसामान्य उपसमूहों का समूचा जालक (अष्टियाँ और प्रतिच्छेद, \cref{exo:b3:representations:6})।
वह अवयवों के क्रम पकड़ने में विफल रहती है: $D_4$ में पाँच अंतर्वलन
हैं और $Q_8$ में एक --- अतः तुल्याकारिता प्रकार \omterm{def:b3:representations:table}{वर्ण सारणी} से सचमुच
सूक्ष्मतर है।
\end{solution}

\begin{solution}{exo:b3:representations:5}
(क) $S_4$ में $(1\,2\,3)$ के \omterm{ex:b3:groups:actions}{केंद्रक} का क्रम $24/8 =
3$ है: वह $\langle(1\,2\,3)\rangle \subseteq A_4$ है। अतः $Z_{A_4}((1\,2\,3))$
का क्रम $3$ है और $A_4$-वर्ग में $12/3
= 4$ अवयव हैं: इसलिए आठों $3$-चक्र
दो $A_4$-वर्गों में बँट जाते हैं (जिनके प्रतिनिधि $(1\,2\,3)$ और उसका
प्रतिलोम हैं)।

(ख) $A_4/V \cong \Z/3\Z$ तीन घात-$1$ \omterm{def:b3:representations:character}{वर्ण} देता है ($\omega = \eu^{2\iu\pi/3}$; $3$-चक्र वर्ग $\bar1, \bar2$
पर जाते हैं); और $\chi_{\mathrm{std}}$ का प्रतिबंध \omterm{def:b3:representations:rep}{अखंडनीय} बना रहता है ($\langle\chi,\chi\rangle = \frac1{12}(9 + 3
\cdot 1 + 0 + 0) = 1$):
\[
\begin{array}{c|cccc}
A_4 & e\,[1] & (1\,2)(3\,4)\,[3] & (1\,2\,3)\,[4] &
(1\,3\,2)\,[4]\\
\hline
\mathbf 1 & 1 & 1 & 1 & 1\\
\chi_\omega & 1 & 1 & \omega & \omega^2\\
\chi_{\bar\omega} & 1 & 1 & \omega^2 & \omega\\
\chi_3 & 3 & -1 & 0 & 0
\end{array}
\]
(ग) अष्टियाँ: $\ker\chi_\omega = \ker\chi_{\bar\omega} = V$; $\ker\chi_3 = \{e\}$ (किसी अन्य प्रविष्टि का मापांक $3$
नहीं है)। \omterm{def:b3:groups:normal}{प्रसामान्य उपसमूह} अष्टियों के प्रतिच्छेद हैं (\cref{exo:b3:representations:6}):
$\{e\}$, $V$, $A_4$ --- विशेष रूप से $A_4$ का क्रम $2$ या सूचकांक
$2$ का कोई \omterm{def:b3:groups:normal}{प्रसामान्य उपसमूह} नहीं है।
\end{solution}

\begin{solution}{exo:b3:representations:6}
(क) यदि $\chi(g) = \chi(e) = n$ हो: $\abs{\chi(g)} \leq n$ में बराबरी से $n\lambda = n$ के साथ $\rho(g) = \lambda\,\mathrm{id}$ (\cref{prop:b3:representations:charbasics})
अनिवार्य है: अतः $\rho(g) = \mathrm{id}$। विलोम स्पष्ट है।

(ख) मान लीजिए $N \trianglelefteq G$। $G/N$ का \omterm{ex:b3:representations:examples}{नियमित निरूपण} निष्ठ है; उसका $G/N$ के अखंडनीयों
में विघटन कीजिए और उन्हें $G$ तक उठाइए (\cref{prop:b3:representations:onedim}(ख)): $G$ के ऐसे
\omterm{def:b3:representations:rep}{अखंडनीय} \omterm{def:b3:representations:character}{वर्ण} $\chi_{i_1}, \dots$ जिनकी अष्टियाँ $N$ को धारण करती हैं और जिनकी
\emph{उभयनिष्ठ} अष्टि ठीक $\{e\}$ का पूर्व-प्रतिबिंब है, अर्थात्
$N$ (भागफल पर निष्ठता)। अतः $N = \bigcap_j \ker\chi_{i_j}$।

(ग) यदि $G$ सरल हो: किसी अतुच्छ \omterm{def:b3:representations:rep}{अखंडनीय} $\chi$ के लिए $\ker\chi \trianglelefteq G$ $G$
नहीं है (समूचे $G$ पर तुच्छ कोई \omterm{def:b3:representations:rep}{अखंडनीय निरूपण} तुच्छ \omterm{def:b3:representations:character}{वर्ण} ही होता
है), अतः $\ker\chi = \{e\}$। विलोमतः मान लीजिए सभी अतुच्छ अष्टियाँ तुच्छ हैं,
और $N \neq
G$ के साथ $N \trianglelefteq G$ लीजिए। (ख) में $N$ को अष्टियों के प्रतिच्छेद
के रूप में व्यक्त करने में कोई न कोई सम्मिलित \omterm{def:b3:representations:character}{वर्ण} अतुच्छ है (यदि
सभी तुच्छ होते, तो प्रतिच्छेद $G$ होता), और उसकी अष्टि $\{e\}$ है:
अतः $N =
\{e\}$। इसलिए एकमात्र \omterm{def:b3:groups:normal}{प्रसामान्य उपसमूह} $\{e\}$ और $G$ हैं।
\end{solution}

\begin{solution}{exo:b3:representations:7}
क्रम $8$ का अनाबेली: घात-$1$ वर्णों की संख्या $[G : D(G)]$ है, जो $8$
का उचित विभाजक है (अनाबेली: $D(G) \neq
\{e\}$), और $\sum n_i^2 = 8$। $k$ एकों और शेष
घातों $\geq 2$ के साथ: वर्गों $\geq 4$ के साथ $8 - k \equiv 0$, और $k
\mid 8$, $k < 8$।
$k = 4$: एक घात $2$ --- संगत। $k =
2$: $6$ बचता है, जो वर्गों $\geq 4$
का योग नहीं है। $k = 1$: असंभव, क्योंकि $k = [G : D(G)] \geq 2$ --- $G/D(G)$ कोई अतुच्छ
आबेली $2$-समूह है, क्योंकि $G$ $2$-समूह है जिसमें $p$-समूहों
की \omterm{def:b3:groups:derived}{विलेयता} (\cref{ex:b3:groups:solvableexamples}) से $D(G) \neq G$। अतः प्रतिरूप $(1,1,1,1,2)$ है। $\chi_5$ की निष्ठता:
चारों घात-$1$ \omterm{def:b3:representations:character}{वर्ण} अपनी अष्टियों में $D(G)$ को धारण करते हैं; और
यदि किसी न्यूनतम प्रसामान्य $N$ के लिए $\ker\chi_5
\supseteq N \neq \{e\}$ हो ($N
\subseteq D(G)$ हो या न
हो --- कोई अतुच्छ $N \subseteq \ker\chi_5$ लीजिए), तो $N$ पाँचों अष्टियों में होता,
जिनका प्रतिच्छेद तुच्छ है (\omterm{ex:b3:representations:examples}{नियमित निरूपण} निष्ठ है): विरोधाभास। क्रम
$p^3$ का अनाबेली: $Z(G)$ का क्रम $p$ है (क्रम $p^2$ से $G/Z$ चक्रीय
और $G$ आबेली हो जाता), क्रम $p^2$ वाला $G/Z(G)$ आबेली है, अतः $D(G) \subseteq Z(G)$,
और $D(G) \ne \{e\}$: इसलिए $D(G) = Z(G)$, जिससे घात $1$ के $p^2$ \omterm{def:b3:representations:character}{वर्ण} मिलते हैं। शेष
घातें $\sum n_i^2 = p^3 -
p^2$ संतुष्ट करती हैं जहाँ प्रत्येक $n_i > 1$ $\abs G$ को विभाजित
करती है (\cref{pb:b3:representations:1}, प्रश्न 8) अतः $n_i \in
\{p\}$ ($n_i = p^2$ से अधिक हो जाता: $p^4 > p^3 - p^2$):
अर्थात् घात $p$ के ठीक $p
- 1$ \omterm{def:b3:representations:character}{वर्ण}।
\end{solution}

\begin{solution}{exo:b3:representations:8}
$V, W$ में $G, H$ के निरूपणों $\rho, \sigma$ के लिए $V \otimes W$ पर $G \times H$ का \omterm{def:b3:representations:rep}{निरूपण}
$\rho \boxtimes \sigma$ परिभाषित कीजिए --- मूर्त रूप में, आव्यूहों पर: $(\rho\boxtimes
\sigma)(g,h)$ क्रोनेकर
गुणनफल $\rho(g)\otimes
\sigma(h)$ है, जिसका संचिह्न $\operatorname{tr}\rho(g)\operatorname{tr}\sigma(h) =
\chi(g)\psi(h)$ है (आव्यूहों $A \otimes B$ के क्रोनेकर
गुणनफल का संचिह्न $\operatorname{tr}A\operatorname{tr}B$ है: उसका विकर्ण $a_{kk}b_{ll}$ है)। अतः $\chi\psi$ कोई
\omterm{def:b3:representations:character}{वर्ण} है, और
\[
\langle\chi\psi, \chi'\psi'\rangle_{G\times H}
= \frac{1}{\abs G\abs H}\sum_{g,h}
\overline{\chi(g)\psi(h)}\,\chi'(g)\psi'(h)
= \langle\chi,\chi'\rangle_G\,\langle\psi,\psi'\rangle_H
= \delta_{\chi\chi'}\delta_{\psi\psi'}.
\]
विशेष रूप से $\langle\chi\psi,\chi\psi\rangle = 1$: अतः प्रत्येक $\chi\psi$ \omterm{def:b3:representations:rep}{अखंडनीय} है (\cref{cor:b3:representations:multiplicity})। ये $r_Gr_H$
भिन्न \omterm{def:b3:representations:rep}{अखंडनीय} \omterm{def:b3:representations:character}{वर्ण} हैं; और $G\times H$ के वर्ग वर्गों के गुणनफल हैं ($(g,h) \sim (g',h')$
घटकशः), अतः उनकी संख्या $r_Gr_H$ है: इसलिए सूची पूर्ण है (\cref{thm:b3:representations:numberirr})। $(\Z/2\Z)^2$
के लिए: चारों चिह्न \omterm{def:b3:representations:character}{वर्ण} $(\pm1)\otimes(\pm1)$ --- अर्थात् \cref{exo:b3:representations:4} की सारणी का ऊपरी-बायाँ
खंड। कोई परिमित आबेली समूह चक्रीय समूहों का गुणनफल होता है (\cref{cor:b3:modules:abelian});
और उसके \omterm{def:b3:representations:rep}{अखंडनीय} \omterm{def:b3:representations:character}{वर्ण} चक्रीय वर्णों के गुणनफल होते हैं (\cref{exo:b3:representations:1}): सभी
घात $1$ के।
\end{solution}

\begin{solution}{exo:b3:representations:9}
(क) $D(A_5) = A_5$ (सरल अनाबेली), अतः एकमात्र घात-$1$ \omterm{def:b3:representations:character}{वर्ण} $\mathbf 1$ है (\cref{prop:b3:representations:onedim})।
हमें ऐसे $n_2^2 + n_3^2
+ n_4^2 + n_5^2 = 59$ चाहिए जिनमें प्रत्येक $n_i \geq 2$; वर्गों $4, 9, 16, 25, 36, 49$ की परीक्षा
करने पर: एकमात्र काम करने वाला बहुसमुच्चय $\{9, 9, 16, 25\}$ है: सबसे बड़े वर्ग
$49$ के साथ शेष $10$ तीन वर्गों $\geq 4$ का योग नहीं है; $36$ के
साथ शेष $23$ भी नहीं ($16 + 4 + 4 = 24$, $9 + 9 + 4 = 22$); सबसे बड़े $25$ के साथ जाँचने
पर $25 + 16 + 9 + 9 = 59$ काम करता है और $25 +
25$, $25 + 16 + 16$, $25 + 16 + 4$ रूपांतर विफल हो जाते
हैं; सबसे बड़े $16$ के साथ: $16\cdot3 = 48 < 59 - 4$। घातें: $1, 3, 3, 4, 5$।

(ख) $5$ बिंदुओं पर क्रमचय: स्थिर बिंदु $(5, 1, 2, 0, 0)$, अतः $\langle\chi_4,\chi_4\rangle = \frac{16 + 0 + 20 + 12 + 12}{60}
= 1$ के साथ
$\chi_4 = (4, 0, 1, -1, -1)$: \omterm{def:b3:representations:rep}{अखंडनीय}। छहों सिलो $5$-उपसमूहों पर क्रिया: कोई अंतर्वलन
ठीक $2$ को स्थिर रखता है (\omterm{ex:b3:groups:actions}{प्रसामान्यक} क्रम $10$ के द्विफलकी हैं,
जिनमें से प्रत्येक में $5$ अंतर्वलन हैं: $15$ अंतर्वलनों के लिए
$30$ आपतन), कोई $3$-अवयव $0$ को स्थिर रखता है ($D_5$ में क्रम
$3$ नहीं), और कोई $5$-अवयव ठीक $1$ को (वह किसी अद्वितीय सिलो
में होता है): अतः क्रमचय \omterm{def:b3:representations:character}{वर्ण} $(6, 2, 0, 1, 1)$, और मानक $\frac{25 + 15 + 20 + 0 + 0}{60} =
1$ के साथ $\chi_5 =
(5, 1, -1, 0, 0)$:
\omterm{def:b3:representations:rep}{अखंडनीय}। दो पंक्तियाँ $(3, a, b, c, d)$, $(3, a', b', c',
d')$ शेष रहती हैं। स्तंभ मानक (\cref{cor:b3:representations:column}):
$(1\,2)(3\,4)$ के वर्ग पर $\abs{Z} = 4$: $1 + 0 + 1 + a^2 + a'^2 = 4$; $e$ के विरुद्ध स्तंभ: $1\cdot1 + 4\cdot 0 + 5\cdot1 + 3(a + a') =
0$: अतः
$a + a' = -2$, $a^2 + a'^2 = 2$: $a = a' = -1$। $3$-चक्रों पर $\abs Z = 3$: $1 + 1 + 1 + b^2 + b'^2 = 3$: $b = b'
= 0$। प्रत्येक
$5$-वर्ग पर $\abs Z = 5$: $e$ के विरुद्ध स्तंभ $1\cdot 1 + 4\cdot(-1) + 5\cdot 0 + 3(c + c') = 0$ पढ़ा जाता है, अतः
$c + c' = 1$; और $c^2 + c'^2 = 3$: $\{c, c'\} = \bigl\{\frac{1+\sqrt5}2,
\frac{1-\sqrt5}2\bigr\}$ --- अर्थात् स्वर्णिम अनुपात और उसका संयुग्म;
दूसरा $5$-वर्ग अदला-बदले हुए मान धारण करता है (दोनों पंक्तियों
को लांबिक होना चाहिए)।

(ग) पूरी हुई सारणी में किसी अतुच्छ पंक्ति की कोई प्रविष्टि प्रथम
स्तंभ के बाहर अपनी घात के बराबर नहीं है: अतः प्रत्येक अष्टि $\{g :
\chi_i(g) = n_i\}$
तुच्छ है। \cref{exo:b3:representations:6}(ग) से $A_5$ सरल है।
\end{solution}

\begin{solution}{exo:b3:representations:10}
$\rho(z)$ प्रत्येक $\rho(g)$ के साथ क्रमविनिमेय है ($z$ केंद्रीय है),
अर्थात् $\rho(z) \in \operatorname{End}_G(V) = \C\,\mathrm{id}$ (शूर): $\rho(z) = \lambda_z\,\mathrm{id}$, और $z \mapsto \lambda_z$ गुणनात्मक है: अतः कोई समाकारिता
$Z(G) \to \C^\times$। तब $\abs{\lambda_z} = 1$ (इकाई-मूल) के साथ $\chi(z) = n\lambda_z$: $\abs{\chi(z)} = n$। यदि $\rho$ निष्ठ हो,
तो $\lambda$ $Z(G)$ पर एकैकी है ($\rho(z) = \mathrm{id} \iff \lambda_z =
1$), अतः $Z(G)$ $\C^\times$ में अंतःस्थापित
हो जाता है; और किसी क्षेत्र के गुणनात्मक समूह का कोई परिमित उपसमूह
चक्रीय होता है (\cref{thm:b3:galois:cyclic})।
\end{solution}

\begin{solution}{exo:b3:representations:11}
(क) समतुल्यवर्तिता: $\rho(h)p_i\rho(h)^{-1}$ योग को पुनःसूचकांकित कर देता है ($g \mapsto hgh^{-1}$,
और $\chi_i$ कोई \omterm{def:b3:representations:character}{वर्ग फलन} है): अतः $p_i$ क्रिया के साथ क्रमविनिमेय है।
\omterm{def:b3:representations:character}{वर्ण} $\chi_j$ वाले किसी \omterm{def:b3:representations:rep}{अखंडनीय} $W$ पर शूर प्रतिबंध को अदिश $\lambda\,\mathrm{id}_W$
बना देते हैं; और संचिह्न लेने पर (प्रथम लांबिकता)
\[
\lambda\,n_j = \frac{n_i}{\abs G}\sum_g
\overline{\chi_i(g)}\,\chi_j(g) = n_i\,\langle\chi_i,
\chi_j\rangle = n_i\,\delta_{ij}
\]
अतः $\lambda = \delta_{ij}$।

(ख) $V$ का अखंडनीयों में विघटन कीजिए (मास्के): $p_i$ \omterm{def:b3:representations:character}{वर्ण} $\chi_i$
वाले योज्य पदों पर तत्समक की भाँति और शेष सब पर $0$ की भाँति क्रिया
करता है, अतः $p_i$ उनके योग $V_i$ पर, शेष के योग के अनुदिश, प्रक्षेप
है; और प्रतिबिंब $V_i$ चुने गए विघटन पर निर्भर नहीं करता (वह $p_i$
से स्थिर सदिशों का समुच्चय है, जो बिना किसी चयन के परिभाषित है)।
$\sum_ip_i$ प्रत्येक \omterm{def:b3:representations:rep}{अखंडनीय} योज्य पद पर तत्समक की भाँति क्रिया करता है:
अतः वह $\mathrm{id}_V$ है। $V_i \cong W_i^{\oplus m_i}$ का सूक्ष्मतर विभाजन $\operatorname{Hom}_G(W_i, V)$ के किसी आधार के
चयन पर निर्भर करता है: वह विहित नहीं है।

(ग) $\varepsilon$ के लिए (घात $1$): $p_\varepsilon =
\frac1{6}\sum_{g}\varepsilon(g)\,\rho(g)$, अर्थात् समूह बीजगणित में
\[
p_\varepsilon = \tfrac16\bigl(e - (1\,2) - (1\,3) - (2\,3)
+ (1\,2\,3) + (1\,3\,2)\bigr) .
\]
वर्ग करने पर: $p_\varepsilon^2$ में $g$ का गुणांक $\frac1{36}\sum_{h}\varepsilon(h)\varepsilon(h^{-1}g) =
\frac1{36}\,\varepsilon(g)\sum_h\varepsilon(h)^2 =
\frac{6}{36}\varepsilon(g)$ है: अतः $p_\varepsilon^2 =
p_\varepsilon$। (\omterm{ex:b3:representations:examples}{नियमित
निरूपण} में उसका प्रतिबिंब $\sum_g\varepsilon(g)e_g$ से फैली रेखा है: अर्थात् चिह्न
\omterm{def:b3:representations:rep}{निरूपण} बहुलता $1$ के साथ आता है, जैसा व्यापक सिद्धांत माँगता है।)
\end{solution}

\begin{solution}{exo:b3:representations:12}
(क) वर्गों को इस क्रम में रखिए: $e$ [1], स्थानांतरण [6],
$3$-चक्र [8], $4$-चक्र [6], द्विक स्थानांतरण [3]। दूसरा रैखिक
\omterm{def:b3:representations:character}{वर्ण} $\chi_2 = \varepsilon$ है, जिसके मान $1, -1, 1, -1, 1$ हैं। घात-$2$ \omterm{def:b3:representations:character}{वर्ण} $\chi_3$ के लिए,
प्रत्येक स्तंभ की तत्समक स्तंभ के साथ स्तंभ लांबिकता ($g \neq e$ के लिए
$\sum_in_i\chi_i(g) = 0$) स्थानांतरणों पर $1 - 1 + 2\chi_3 + 3(\chi_4 +
\chi_5) = 0$ देती है; और चिह्न-युक्ति $\chi_5 = \varepsilon\chi_4$ (कोई
घात-$3$ \omterm{def:b3:representations:character}{वर्ण} गुणा कोई रैखिक \omterm{def:b3:representations:character}{वर्ण} फिर से \omterm{def:b3:representations:rep}{अखंडनीय} होता है --- वही
मानक) $\chi_4 + \chi_5$ को विषम वर्गों पर लुप्त कर देती है: वहाँ $\chi_3 = 0$। $3$-चक्रों
पर: सम वर्गों पर $\chi_5 = \chi_4$ के साथ $1 + 1 +
2\chi_3(c_3) + 3(\chi_4 + \chi_5)(c_3) = 0$; $c_3$ के स्तंभ की स्वयं के साथ
तुलना $\abs{\chi_3}^2$ आँकड़ा देती है; और छोटे निकाय को हल करने पर ($\chi_3$ की
$\mathbf 1$ तथा $\varepsilon$ के साथ पंक्ति लांबिकता का भी उपयोग कीजिए): $\chi_3 = (2, 0, -1, 0, 2)$,
फिर $\chi_4 =
(3, 1, 0, -1, -1)$ और $\chi_5 = \varepsilon\chi_4 = (3, -1,
0, 1, -1)$। पूरी सारणी:
\begin{center}
\begin{tabular}{c|ccccc}
 & $e$ & $6\,t$ & $8\,c_3$ & $6\,c_4$ & $3\,v$\\
\hline
$\chi_1$ & $1$ & $1$ & $1$ & $1$ & $1$\\
$\chi_2$ & $1$ & $-1$ & $1$ & $-1$ & $1$\\
$\chi_3$ & $2$ & $0$ & $-1$ & $0$ & $2$\\
$\chi_4$ & $3$ & $1$ & $0$ & $-1$ & $-1$\\
$\chi_5$ & $3$ & $-1$ & $0$ & $1$ & $-1$\\
\end{tabular}
\end{center}
(सभी पंक्तियों का मानक $1$ है; सभी स्तंभ लांबिक हैं: जाँच पूरी।)

(ख) इन घातों के साथ वर्ग आँकड़ा $1, 6, 8, 6, 3$ $S_4$ का ही है (चक्र प्रकार
$e$, $2$, $3$, $4$, $2{+}2$)। अष्टियाँ: $\ker\chi_2 = \{g : \varepsilon(g) = 1\} = A_4$ (वर्ग $e, c_3, v$:
$1 + 8 + 3 = 12$, सूचकांक $2$); $\ker\chi_3 = \{g : \chi_3(g) = 2\}$ = वर्ग $e, v$: क्रम $4$ का क्लाइन
समूह $V_4$, प्रसामान्य। $\chi_4, \chi_5$ निष्ठ हैं ($\chi_i(g) = n_i$ केवल $e$ पर)।
अष्टियों के प्रतिच्छेद: $\{e\}$, $V_4$, $A_4$, $G$ --- और \cref{exo:b3:representations:6}(ख)
से ये $S_4$ के \emph{सभी} \omterm{def:b3:groups:normal}{प्रसामान्य उपसमूह} हैं।

(ग) $V_4 \subseteq \ker$ वाले \omterm{def:b3:representations:character}{वर्ण} $\chi_1,
\chi_2, \chi_3$ हैं: वे $G/V_4$ के माध्यम से गुणनखंडित होते
हैं, जो क्रम $6$ का समूह है और जिसकी \omterm{def:b3:representations:rep}{अखंडनीय} घातें $1, 1, 2$ हैं ---
अर्थात् $S_3$ की सारणी। चूँकि भागफल की सारणी क्रम $6$ के दोनों
समूहों के बीच एक पूर्ण निश्चर \emph{है} ($\Z/6\Z$ के छह रैखिक \omterm{def:b3:representations:character}{वर्ण}
होते), अतः $G/V_4 \cong
S_3$: भागफल सारणी के भीतर उन पंक्तियों के खंड के रूप
में दिख जाता है जिनकी अष्टि में $V_4$ है।
\end{solution}

\begin{solution}{pb:b3:representations:1}
\textbf{1.} यदि $\alpha^n + c_{n-1}\alpha^{n-1} + \dots + c_0 =
0$ ($c_i \in \Z$) हो, तो $\alpha^n \in \Z\text{-विस्तार}(1, \dots,
\alpha^{n-1})$, और आगमन से प्रत्येक घात
भी: अतः $\Z[\alpha]$ $1, \alpha, \dots, \alpha^{n-1}$ से जनित है। विलोमतः मान लीजिए $\Z[\alpha] = \Z g_1 + \dots + \Z g_m$। $M = (m_{ij}) \in M_m(\Z)$ के
साथ $\alpha g_i =
\sum_j m_{ij}g_j$ लिखिए: सदिश $g = (g_i)$ $(\alpha I - M)g = 0$ संतुष्ट करता है; सहगुणक आव्यूह
से गुणा करने पर प्रत्येक $i$ के लिए $\det(\alpha I - M)\,g_i = 0$, और चूँकि $1 \in \Z[\alpha]$ $g_i$
का कोई $\Z$-संचय है, अतः $\det(\alpha I - M) = 0$: इसलिए $\alpha$ मोनिक $\det(XI - M) \in \Z[X]$ का मूल है।

\textbf{2.} यदि $\Z[\alpha]$ $n-1$ तक $\alpha$ की घातों से और $\Z[\beta]$ $m - 1$ तक
$\beta$ की घातों से फैला हो, तो $\Z[\alpha,\beta]$ $nm$ गुणनफलों $\alpha^i\beta^j$ से फैला है
(किसी भी एकपदी को घटाइए)। उपवलय $\Z[\alpha + \beta]$ और $\Z[\alpha\beta]$ \omterm{def:b3:modules:free}{परिमिततः जनित}
$\Z$-मॉड्यूल $\Z[\alpha, \beta]$ के $\Z$-उपमॉड्यूल हैं, अतः \omterm{def:b3:modules:free}{परिमिततः जनित} ($\Z[\alpha,\beta]$,
जो $nm$ अवयवों से जनित है, $\Z^{nm}$ का प्रतिबिंब है; और कोई उपमॉड्यूल
$\Z^{nm}$ के किसी उपमॉड्यूल पर पीछे खिंच आता है, जो \cref{thm:b3:modules:submodule} से कोटि $\leq nm$
का मुक्त है, और उसका प्रतिबिंब जनित करता है)। प्रश्न 1 से $\alpha + \beta$ और
$\alpha\beta$ \omterm{def:b3:galois:algebraic}{बीजीय} पूर्णांक हैं।

\textbf{3.} यदि $\frac pq$ (निम्नतम पदों में) घात $n$ के किसी मोनिक
पूर्णांक बहुपद का मूल हो, तो परिमेय मूल प्रमेय (हर हटाइए: $p^n = -q(\cdots)$)
$q \mid p^n$ दे देती है, अतः $q = \pm1$।

\textbf{4.} $\chi(g)$ इकाई-मूलों का योग है (\cref{prop:b3:representations:charbasics}), और उनमें से प्रत्येक
\omterm{def:b3:galois:algebraic}{बीजीय} पूर्णांक है ($X^N - 1$ का मूल); अतः प्रश्न 2 से निष्कर्ष निकलता
है।

\textbf{5.} $h \in G$ के लिए: $\rho(h)S_C\rho(h)^{-1} = \sum_{g\in
C}\rho(hgh^{-1}) = S_C$ ($C$ कोई वर्ग है)। शूर से $S_C =
\omega_C\,\mathrm{id}$;
और संचिह्न लेने पर $\abs C\,\chi(g_C) =
\omega_C\, n$।

\textbf{6.} $a(z) = \#\{(x,y) \in C\times
C' : xy = z\}$ के साथ $S_CS_{C'} = \sum_{x \in C, y \in C'}\rho(xy) =
\sum_{z \in G} a(z)\rho(z)$। $h$ से संयुग्मन $z$ के हलों को
$hzh^{-1}$ के हलों के साथ एकैकी रूप से जोड़ देता है: अतः $a$ $\N$ में
मान लेने वाला कोई \omterm{def:b3:representations:character}{वर्ग फलन} है, इसलिए $S_CS_{C'} = \sum_{C''}a_{CC'C''}S_{C''}$। सर्वत्र $S_C = \omega_C\,\mathrm{id}$ प्रतिस्थापित
करके और अदिशों की पहचान करके: $\omega_C\omega_{C'} = \sum_{C''}a_{CC'C''}\,\omega_{C''}$।

\textbf{7.} मान लीजिए $M$ $1$ और सभी गुणनफलों $\omega_{C_1}\cdots\omega_{C_k}$ से फैला
$\Z$-मॉड्यूल है; प्रश्न 6 से ऐसा प्रत्येक गुणनफल $1$ और $\omega_C$
के किसी $\Z$-संचय तक सिमट जाता है: अतः $M$ \omterm{def:b3:modules:free}{परिमिततः जनित} है, और
प्रत्येक $C$ के लिए $\omega_C M
\subseteq M$। विशेष रूप से $\Z[\omega_C]
\subseteq M$ \omterm{def:b3:modules:free}{परिमिततः जनित} है
(उपमॉड्यूल, प्रश्न 2 की भाँति), और प्रश्न 1 $\omega_C$ को \omterm{def:b3:galois:algebraic}{बीजीय} पूर्णांक
बना देता है।

\textbf{8.} घात $n$ के किसी \omterm{def:b3:representations:rep}{अखंडनीय} $\chi$ के लिए:
\[
\frac{\abs G}{n} = \frac{\abs G}{n}\langle\chi,\chi\rangle
= \frac1n \sum_{g}\chi(g)\overline{\chi(g)}
= \sum_{C}\frac{\abs C\,\chi(g_C)}{n}\,\overline{\chi(g_C)}
= \sum_C \omega_C\,\overline{\chi(g_C)} .
\]
प्रत्येक $\overline{\chi(g_C)} = \chi(g_C^{-1})$ \omterm{def:b3:galois:algebraic}{बीजीय} पूर्णांक है (प्रश्न 4), अतः दायाँ पक्ष भी
(प्रश्न 2, 7); और वह परिमेय है, इसलिए पूर्णांक (प्रश्न 3): $n \mid \abs
G$।

\textbf{9.} बेज़ू: $u, v \in \Z$ के साथ $u\abs C + vn = 1$। तब
\[
\frac{\chi(g_C)}{n} = u\,\frac{\abs C\,\chi(g_C)}{n} +
v\,\chi(g_C) = u\,\omega_C + v\,\chi(g_C),
\]
जो कोई \omterm{def:b3:galois:algebraic}{बीजीय} पूर्णांक है।

\textbf{10.} मान लीजिए $0 < \abs{\chi(g_C)} < n$ और $\alpha
= \chi(g_C)/n$ लीजिए। सभी मान $\chi(g)$ $\Q(\zeta_N)$,
$N =
\abs G$ में हैं ($N$-वें इकाई-मूलों के योग)। $\sigma \in
\operatorname{Gal}(\Q(\zeta_N)/\Q)$ के लिए: $\sigma$
इकाई-मूलों को इकाई-मूलों पर भेजता है ($\sigma(\zeta^k) = \zeta^{ak}$, \cref{thm:b3:galois:cyclotomicirred}), अतः $\sigma(\chi(g_C))$ फिर
से $n$ इकाई-मूलों का योग है: $\abs{\sigma(\alpha)}
\leq 1$; साथ ही $\sigma(\alpha)$ \omterm{def:b3:galois:algebraic}{बीजीय} पूर्णांक
है (उसका \omterm{def:b3:galois:algebraic}{न्यूनतम बहुपद} $\alpha$ के ही समान है)। गुणनफल $P =
\prod_\sigma \sigma(\alpha)$ समूचे
\omterm{def:b3:galois:galois}{गाल्वा समूह} से स्थिर रहता है, अतः परिमेय (\cref{thm:b3:galois:fundamental}(1)), और वह
\[
0 < \abs P \leq \abs\alpha < 1
\]
वाला \omterm{def:b3:galois:algebraic}{बीजीय} पूर्णांक है (कोई गुणनखंड लुप्त नहीं होता: $\sigma(\alpha) = 0$ से $\alpha =
0$
अनिवार्य हो जाता)। यह प्रश्न 3 का खंडन करता है। अतः $\chi(g_C) = 0$ या $\abs{\chi(g_C)} = n$,
और पिछली स्थिति में $\rho(g_C)$ अदिश है (\cref{prop:b3:representations:charbasics})।

\textbf{11.} स्तंभ लांबिकता ($C \neq \{e\}$): $\sum_i \chi_i(e)\overline{\chi_i(g_C)} = 0$, अर्थात् $1 +
\sum_{i \geq 2} n_i\overline{\chi_i(g_C)} = 0$। यदि
$p \nmid n_i$ वाला प्रत्येक अतुच्छ $\chi_i$ $g_C$ पर लुप्त हो जाता, तो शेष को
उनके गुणनखंड $p$ के अनुसार समूहित करने पर
\[
-\frac1p = \sum_{i \geq 2,\ p \mid n_i}
\frac{n_i}{p}\,\overline{\chi_i(g_C)},
\]
कोई \omterm{def:b3:galois:algebraic}{बीजीय} पूर्णांक होता --- जो प्रश्न 3 का खंडन करता है। अतः किसी
अतुच्छ $\chi_i$ के लिए $p \nmid n_i$ और $\chi_i(g_C) \neq 0$; और चूँकि $\abs C = p^k$, अतः $\gcd(\abs C, n_i) = 1$, और
प्रश्न 10 $\rho_i(g_C)$ को अदिश बना देता है। अब $G$ सरल अनाबेली: $\chi_i$
निष्ठ है (\cref{exo:b3:representations:6}(ग)), और $Z_i = \{g : \rho_i(g) \text{ अदिश}\}$ $g_C \neq e$ को धारण करने वाला कोई \omterm{def:b3:groups:normal}{प्रसामान्य
उपसमूह} है ($\rho_i$ के अंतर्गत अदिशों का पूर्व-प्रतिबिंब, और अदिश
प्रतिबिंब का प्रसामान्य --- वस्तुतः केंद्रीय --- उपसमूह बनाते हैं):
अतः $Z_i = G$। तब $\rho_i(G)$ आबेली और निष्ठ है, जिससे $G$ आबेली हो जाता
है: विरोधाभास। \emph{किसी सरल अनाबेली समूह का अभाज्य-घात आकार $> 1$
वाला कोई \omterm{ex:b3:groups:actions}{संयुग्मन वर्ग} नहीं होता।}

\textbf{12.} मान लीजिए $G$ क्रम $p^aq^b$ का सरल है। यदि $b = 0$ हो
($G$ कोई $p$-समूह): $Z(G) \neq \{e\}$ (\cref{thm:b3:groups:pfixed}) प्रसामान्य है, अतः $Z(G) = G$:
आबेली। अन्यथा कोई सिलो $q$-उपसमूह $Q$ और $z \in Z(Q)
\setminus\{e\}$ लीजिए (पुनः
\cref{thm:b3:groups:pfixed})। तब $Q
\subseteq Z_G(z)$, अतः $z$ के वर्ग का आकार $[G : Z_G(z)]$ है जो $[G : Q] = p^a$ को
विभाजित करता है: अर्थात् $p$ की कोई घात। यदि आकार $1$ हो, तो
$z \in Z(G)$: \omterm{ex:b3:groups:actions}{केंद्र} कोई अतुच्छ \omterm{def:b3:groups:normal}{प्रसामान्य उपसमूह} है, अतः $Z(G) = G$, आबेली।
यदि आकार $p^k > 1$ हो: प्रश्न 11 सरल अनाबेली $G$ के लिए उसे मना कर
देता है। किसी भी दशा में क्रम $p^aq^b$ का कोई \omterm{def:b3:groups:simple}{सरल समूह} आबेली है ($\cong \Z/p\Z$)।

\textbf{13.} $\abs G$ पर आगमन ($\abs G = 1$: \omterm{def:b3:groups:derived}{विलेय})। यदि $G$ सरल हो, तो
प्रश्न 12 उसे आबेली बना देता है, अतः \omterm{def:b3:groups:derived}{विलेय}। अन्यथा कोई अतुच्छ उचित
\omterm{def:b3:groups:normal}{प्रसामान्य उपसमूह} $N$ चुनिए: $\abs
N$ और $\abs{G/N}$ फिर से $p^{a'}q^{b'}$ रूप के
और छोटे हैं, अतः आगमन से $N$ और $G/N$ \omterm{def:b3:groups:derived}{विलेय} हैं, और $G$ \omterm{def:b3:groups:derived}{विलेय}
है (\cref{prop:b3:groups:derived})। --- तीन अभाज्यों के साथ मुख्य चरण विफल हो जाता है: किसी
\omterm{def:b3:groups:sylow}{सिलो उपसमूह} का सूचकांक अब अभाज्य घात नहीं रहता, अतः किसी सिलो के
केंद्रीय अवयव के वर्ग का आकार अभाज्य घात होना आवश्यक नहीं। और कोई
न कोई विफलता अनिवार्य है: क्रम $2^2\cdot3\cdot5$ का $A_5$ सरल है और \omterm{def:b3:groups:derived}{विलेय} नहीं।

\textbf{14.} वर्ग और आकार \cref{exo:b3:groups:11}(क) हैं। $c$ के $S_5$-वर्ग का आकार
$24$ है; और यदि वह एक ही $A_5$-वर्ग रहता, तो कक्षा--स्थायीकारी गणना
$\abs{Z_{A_5}(c)} = 60/24$ देती, जो पूर्णांक नहीं है --- मूर्त रूप में, $Z_{S_5}(c) = \langle c\rangle$ का क्रम
$5$ है और वह $A_5$ के भीतर बैठता है, अतः $c$ के $A_5$-वर्ग का
आकार $60/5 =
12$ है: इसलिए $S_5$-वर्ग दो में बँट जाता है ($c$ और $c^2$
$S_5$ में केवल किसी विषम क्रमचय से संयुग्म हैं)।

\textbf{15.} एक रैखिक \omterm{def:b3:representations:character}{वर्ण}: कोई घात-$1$ \omterm{def:b3:representations:rep}{निरूपण} $G/D(G)$ के माध्यम
से गुणनखंडित होता है, और $D(A_5) = A_5$ ($A_5$ सरल अनाबेली है, और $D(A_5)$
प्रसामान्य तथा अतुच्छ है --- $A_5$ आबेली नहीं है)। अतः $n_1 = 1$ और
$n_2^2 + n_3^2 + n_4^2
+ n_5^2 = 59$, जिसमें प्रत्येक $n_i \geq 2$। उपलब्ध वर्ग: $4, 9,
16, 25, 36, 49$। उनमें से चार का
योग $59$ के बराबर: सबसे बड़ा $25$ होना चाहिए ($36 + 4 + 4 + 9 = 53 < 59$ समायोजित नहीं
हो पाता: $36 + 16 + 4 + 4 = 60$, $36 + 9 + 9 + 4 = 58$, $36 +
16 + 9 + 4 = 65$ --- $36$ या $49$ के साथ कोई संयोजन
काम नहीं करता), और $59 - 25 = 34 = 16 + 9 + 9$ (एकमात्र उपाय: $16 + 16 + 4 =
36$, $25 + 9 + 4 = 38$, $25 + 4 + 4 = 33$): अतः
घातें $1, 3, 3,
4, 5$।

\textbf{16.} $\langle\pi, \mathbf 1\rangle = \frac1{60}(5 +
15\cdot1 + 20\cdot2 + 0 + 0) = 1$ (एक \omterm{def:b3:groups:action}{कक्षा} --- बर्नसाइड की गणना), और $\langle\pi, \pi\rangle = \frac1{60}(25 + 15 + 80 +
0 + 0) = 2$: अतः
$\pi$ तुच्छ \omterm{def:b3:representations:character}{वर्ण} को एक बार धारण करता है, और उसका दूसरा घटक कोई एकल
\omterm{def:b3:representations:rep}{अखंडनीय} है। इसलिए $\chi_4 = \pi
- \mathbf 1$ \omterm{def:b3:representations:rep}{अखंडनीय} है, जिसकी घात $4$ और मान $4, 0,
1, -1, -1$
हैं।

\textbf{17.} युग्मों पर कोई अवयव $\{i,j\}$ को स्थिर रखता है तभी और केवल
तभी जब वह $i, j$ को स्थिर रखे या उनकी अदला-बदली करे: गणनाएँ $10$
($e$), $2$ ($t =
(1\,2)(3\,4)$ $\{1,2\}, \{3,4\}$ को स्थिर रखता है), $1$ ($(1\,2\,3)$ $\{4,5\}$
को स्थिर रखता है), $0, 0$ हैं। तब $\langle\pi_{10},
\pi_{10}\rangle = \frac1{60}(100 + 15\cdot4 + 20\cdot1) = 3$: तीन \omterm{def:b3:representations:rep}{अखंडनीय} घटक, प्रत्येक
एक बार। और $\langle\pi_{10}, \mathbf 1\rangle = \frac1{60}(10 + 30 + 20)
= 1$, $\langle\pi_{10}, \chi_4\rangle = \frac1{60}(40 + 0 +
20\cdot1\cdot1 + 0 + 0) = 1$: अतः $\pi_{10} = \mathbf 1 + \chi_4
+ \chi$, जहाँ $\chi$ घात
$10 - 1 - 4 = 5$ और मान $\chi_5 = \pi_{10} - \mathbf 1 - \chi_4 = (5, 1, -1, 0,
0)$ वाला \omterm{def:b3:representations:rep}{अखंडनीय} है।

\textbf{18.} $t$ पर $\chi_2, \chi_3$ के मानों के लिए $a, a'$ लिखिए, और $s = (1\,2\,3)$
पर $b, b'$; चारों वास्तविक हैं ($t$ और $s$ अपने प्रतिलोमों के
संयुग्म हैं)। स्तंभ $e$ के विरुद्ध स्तंभ $t$: $1 + 3a + 3a' + 4\cdot0 + 5\cdot1 = 0$, अतः $a + a'
= -2$;
स्तंभ $t$ स्वयं के साथ: $1 + a^2 + a'^2 + 0 + 1 =
\frac{60}{15} = 4$, अतः $a^2 + a'^2 = 2$; इसलिए $(a + a')^2 = 4
= 2 + 2aa'$ से $aa' = 1$
और $a = a' = -1$ मिलते हैं। स्तंभ $s$ $e$ के विरुद्ध: $1 + 3(b + b') + 4 - 5 = 0$ से $b + b' = 0$
मिलता है; स्तंभ $s$ स्वयं के साथ: $1 + b^2 + b'^2 + 1 + 1 = \frac{60}
{20} = 3$ से $b = b' = 0$।

\textbf{19.} स्तंभ $c$ स्तंभ $e$ के विरुद्ध: $1 + 3(x + y) +
4(-1) + 5\cdot0 = 0$, अतः $x + y = 1$।
स्तंभ $c$ स्तंभ $c^2$ के विरुद्ध (भिन्न वर्ग, लांबिक): $1 + xy + yx + 1
+ 0 = 0$, अतः
$xy = -1$। इस प्रकार $x, y$ $T^2 - T - 1 = 0$ हल करते हैं: $\varphi = \frac{1 +
\sqrt5}2$ के साथ $\{x, y\} = \{\varphi, \bar\varphi\}$। संगति:
$x^2 + y^2 = (x+y)^2 - 2xy = 3 =
\frac{60}{12} - 2$ --- जो स्व-स्तंभ सर्वसमिका $1 +
x^2 + y^2 + 1 + 0 = 5$ से मेल खाता है। सारणी:
\begin{center}
\begin{tabular}{c|ccccc}
 & $e$ & $15\,t$ & $20\,s$ & $12\,c$ & $12\,c^2$\\
\hline
$\chi_1$ & $1$ & $1$ & $1$ & $1$ & $1$\\
$\chi_2$ & $3$ & $-1$ & $0$ & $\varphi$ & $\bar\varphi$\\
$\chi_3$ & $3$ & $-1$ & $0$ & $\bar\varphi$ & $\varphi$\\
$\chi_4$ & $4$ & $0$ & $1$ & $-1$ & $-1$\\
$\chi_5$ & $5$ & $1$ & $-1$ & $0$ & $0$\\
\end{tabular}
\end{center}

\textbf{20.} $\varphi^2 + \bar\varphi^2 = (\varphi +
\bar\varphi)^2 - 2\varphi\bar\varphi = 1 + 2 = 3$ का उपयोग करते हुए $\norm{\chi_2}^2 = \frac1{60}\bigl(9 + 15\cdot1 +
0 + 12\varphi^2 + 12\bar\varphi^2\bigr) = \frac{9 + 15 +
36}{60} = 1$। इसी प्रकार $\langle\chi_2, \chi_3\rangle = \frac1{60}(9 + 15 + 0 +
12(2\varphi\bar\varphi)) = \frac{9 + 15 - 24}{60} = 0$।
\omterm{thm:b3:galois:finitefields}{फ्रोबेनियस} विभाज्यता: $1, 3, 3, 4, 5$ सभी $60$ को विभाजित करते हैं। अपरिमेय
मान $\varphi, \bar\varphi$ $S_5$ के साथ दिखने वाला अंतर हैं: $S_5$ में प्रत्येक अवयव
अपने चक्रीय समूह के उन सभी जनकों का संयुग्म होता है जिनका चक्र प्रकार
वही हो --- विशेष रूप से $c \sim c^2$ --- जिससे परिमेय (वस्तुतः पूर्णांक)
\omterm{def:b3:representations:character}{वर्ण} मान अनिवार्य हो जाते हैं; और $A_5$ में पाँच-चक्रों का विभाजन
$\Q(\sqrt5)$ के लिए द्वार खोल देता है।

\textbf{21.} कोई \omterm{def:b3:groups:normal}{प्रसामान्य उपसमूह} $N$ वर्गों का संघ होता है,
$e$ को धारण करता है, और $\abs N \mid 60$। संभावित योग: $1{+}15 = 16$, $1{+}20 = 21$, $1{+}12 = 13$,
$1{+}24 = 25$, $1{+}15{+}20 = 36$, $1{+}15{+}12 = 28$, $1{+}15{+}24 = 40$, $1{+}20{+}12 = 33$, $1{+}20{+}24 = 45$, $1{+}12{+}12 = 25$, $1{+}15{+}20{+}12 = 48$, $1{+}15{+}20{+}24 = 60 - 12 = \dots$ ---
$1$ को धारण करने वाले सभी उचित उप-योग सूचीबद्ध करने पर: $13, 16,
21, 25, 28, 33, 36, 40, 45, 48, 13{+}\dots$ में
से कोई भी $60$ को विभाजित नहीं करता, सिवाय $1$ के: अतः $N = \{e\}$ या
$A_5$। सरलता, पाँच संख्याओं से पढ़ ली गई। और प्रश्न 11 की कसौटी भी
दिख जाती है: वर्ग-आकार $15 = 3\cdot5$, $20 = 4\cdot5$, $12 = 4\cdot3$ सभी दो अभाज्यों के संयुक्त
हैं --- कोई अभाज्य-घात वर्ग नहीं, ठीक वैसे ही जैसे बर्नसाइड की कसौटी
किसी \omterm{def:b3:groups:simple}{सरल समूह} से माँगती है।

\textbf{22.} कोण $\theta$ से $\R^3$ के किसी घूर्णन के अभिलक्षणिक मान
$1, \eu^{\iu\theta}, \eu^{-\iu\theta}$ हैं: संचिह्न $1 +
2\cos\theta$। $\theta = \frac{2\pi}5$ के लिए: $2\cos\frac{2\pi}5
= \frac{\sqrt5 - 1}2$, अतः संचिह्न $1 + \frac{\sqrt5-1}2 =
\frac{1+\sqrt5}2 = \varphi$
है। $\operatorname{Gal}(\Q(\sqrt5)/\Q)$ का अतुच्छ अवयव $\tau$ किसी \omterm{def:b3:representations:table}{वर्ण सारणी} पर प्रविष्टिशः लगाया
जाए तो वर्णों को वर्णों पर भेजता है (वह परिभाषक बीजगणित के साथ
क्रमविनिमेय है: $\tau\circ\chi$ उस \omterm{def:b3:representations:rep}{निरूपण} का \omterm{def:b3:representations:character}{वर्ण} है जो आव्यूहों को प्रविष्टियों
पर $\tau$ के माध्यम से ले जाकर प्राप्त होता है; अथवा सारगर्भित रूप
से: लांबिकता संबंध $\Q$-परिमेय हैं, अतः $\tau$ उनके हलों का क्रमचय
करता है); $\tau$ $\chi_1, \chi_4, \chi_5$ को स्थिर रखता है (परिमेय मान) और इसलिए उसे
$\chi_2$ तथा $\chi_3$ की अदला-बदली करनी ही होगी: अर्थात् दूसरा घात-$3$
\omterm{def:b3:representations:character}{वर्ण} संयुग्मित मान धारण करता है, और कोई आव्यूह परिकलित नहीं करना
पड़ा। ज्यामितीय रूप से, दोनों \omterm{def:b3:representations:rep}{निरूपण} बीसफलकीय क्रिया और उसका $A_5$
की किसी बाह्य स्वसमाकारिता (किसी स्थानांतरण से संयुग्मन) के साथ
संयोजन हैं, जो पाँच-चक्रों के दोनों वर्गों की अदला-बदली कर देता है।

\textbf{23.} $z \in Z(G)$ के लिए $\rho(z)$ प्रत्येक $\rho(g)$ के साथ क्रमविनिमेय
है, अतः शूर की प्रमेयिका से $\rho(z) = \lambda\,
\mathrm{id}$; और चूँकि $z$ का क्रम परिमित
है, $\lambda$ कोई इकाई-मूल है, तथा $\abs{\chi(z)} = \abs\lambda\,n = n$। तब
\[
\abs G = \abs G\,\langle\chi, \chi\rangle
= \sum_{g \in G}\abs{\chi(g)}^2
\geq \sum_{z \in Z(G)}\abs{\chi(z)}^2
= \abs{Z(G)}\,n^2,
\]
अर्थात् $n^2 \leq [G : Z(G)]$। क्रम $8$ के किसी अनाबेली समूह ($D_4$ या $Q_8$) के
पाँच \omterm{ex:b3:groups:actions}{संयुग्मन वर्ग} होते हैं, अतः $\sum n_i^2 = 8$ वाली पाँच \omterm{def:b3:representations:rep}{अखंडनीय} घातें; और
$8$ को $\geq 1$ जैसे पाँच वर्गों के योग के रूप में लिखने का एकमात्र
उपाय $1 + 1 + 1 + 1
+ 4$ है: अतः घातें $1, 1, 1, 1, 2$। दोनों समूहों के \omterm{ex:b3:groups:actions}{केंद्र} का क्रम $2$
है, और घात-$2$ \omterm{def:b3:representations:character}{वर्ण} बराबरी प्राप्त कर लेता है: $2^2 = 4 = [G : Z(G)]$ --- अर्थात्
परिबंध तीक्ष्ण है। $A_5$ के लिए $Z = \{e\}$ और परिबंध $n^2 \leq 60$ पढ़ा जाता है:
जिसे $1, 3, 3, 4, 5$ खुली जगह के साथ संतुष्ट करता है ($25 \leq 60$), और ऐसा होना ही
चाहिए, क्योंकि बराबरी से (उसी शृंखला द्वारा) $\chi$ को \omterm{ex:b3:groups:actions}{केंद्र} के
बाहर लुप्त होना पड़ता।

\textbf{24.} $g$ के क्रम $m$ के लिए $\rho(g)^m = \mathrm{id}$, अतः $\rho(g)$ $X^m - 1$ से
नष्ट हो जाता है, जो $\C$ पर सरल मूलों के साथ विभाजित है: इसलिए वह
विकर्णीय है, जिसके अभिलक्षणिक मान $\lambda_1, \dots, \lambda_n$ इकाई-मूल हैं, और अभिलक्षणिक
आधार $(e_i)$ है। गुणनफल $e_ie_j$ ($i \leq j$) अभिलक्षणिक मानों $\lambda_i\lambda_j$ के साथ
$\operatorname{Sym}^2V$ का अभिलक्षणिक आधार बनाते हैं, और $e_i \wedge e_j$ ($i < j$) $\Lambda^2V$ का; और
चूँकि
\[
\sum_{i<j}\lambda_i\lambda_j
= \frac{(\sum_i\lambda_i)^2 - \sum_i\lambda_i^2}2
= \frac{\chi(g)^2 - \chi(g^2)}2,
\]
तथा सममित योग $\sum_i\lambda_i^2$ को घटाने के बजाय जोड़ता है, अतः दोनों सूत्र
निकल आते हैं। वर्गों $(e, (2,2)\text{-},
3\text{-चक्र}, c, c^2)$ पर $\chi_2 = (3, -1,
0, \varphi, \bar\varphi)$ के लिए: वर्ग करना द्विक स्थानांतरणों
को $e$ पर, तीन-चक्रों को तीन-चक्रों पर, और $c$ के वर्ग को $c^2$
के वर्ग पर तथा विलोमतः भेजता है ($A_5$ में $c^{-1} \sim
c$, अतः $c^4 \sim c$)। अतः
$\chi_2(g^2)$ $(3,
3, 0, \bar\varphi, \varphi)$ के रूप में पढ़ा जाता है, और $\varphi^2 = \varphi
+ 1$, $\bar\varphi = 1 - \varphi$ का उपयोग करते
हुए:
\[
\Lambda^2\chi_2 = (3, -1, 0, \varphi, \bar\varphi) = \chi_2,
\qquad
\operatorname{Sym}^2\chi_2 = (6, 2, 0, 1, 1) .
\]
वर्ग-आकार $1, 15, 20, 12,
12$ के साथ पिछले का विघटन करने पर: $\langle\cdot, \mathbf 1\rangle = \frac1{60}(6 +
15\cdot2 + 0 + 12 + 12) = 1$; $\langle\cdot, \chi_5\rangle =
\frac1{60}(30 + 30) = 1$; $\langle\cdot, \chi_4\rangle =
\frac1{60}(24 - 12 - 12) = 0$;
$\langle\cdot, \chi_2\rangle =
\frac1{60}(18 - 30 + 12(\varphi + \bar\varphi)) = 0$, और इसी प्रकार $\chi_3$ के लिए। अतः $\operatorname{Sym}^2\chi_2 =
\mathbf 1 + \chi_5$ (विमाएँ $6 = 1 + 5$) और
$\chi_2\otimes\chi_2 = \mathbf 1 + \chi_2 + \chi_5$
(विमाएँ $9 = 1 + 3 + 5$)। ज्यामिति: समतुल्यवर्ती तुल्याकारिता $\Lambda^2\R^3 \to \R^3$, $u \wedge v \mapsto u
\times v$,
किसी घूर्णन समूह के लिए ठीक $\Lambda^2\chi_2 = \chi_2$ है; सममित वर्ग का योज्य पद $\mathbf 1$
निश्चर द्विघात रूप $x^2 + y^2 + z^2$ है, और $\chi_5$ संचिह्न-रहित सममित प्रदिशों
(प्रसंवादी द्विघातों) की पाँच-विमीय समष्टि पर रहता है।

\textbf{25.} $c^2$ के स्तंभ के विरुद्ध $c$ का स्तंभ:
\[
1\cdot1 + \varphi\bar\varphi + \bar\varphi\varphi +
(-1)(-1) + 0 = 1 - 1 - 1 + 1 + 0 = 0,
\]
जैसा भिन्न वर्गों के लिए लांबिकता माँगती है ($\varphi
\bar\varphi = -1$)। $c$ का स्तंभ
स्वयं के विरुद्ध: $1 +
\varphi^2 + \bar\varphi^2 + 1 + 0 = 1 + 3 + 1 = 5 = 60/12 =
\abs{Z_{A_5}(c)}$। अतत्समक चारों स्तंभों पर नियमित \omterm{def:b3:representations:character}{वर्ण} $\sum_in_i\chi_i$:
\[
1 - 3 - 3 + 0 + 5 = 0, \qquad
1 + 0 + 0 + 4 - 5 = 0,
\]
द्विक स्थानांतरणों और तीन-चक्रों पर, तथा $c$ के वर्ग पर ($c^2$
का वर्ग उसका \omterm{def:b3:galois:galois}{गाल्वा} संयुग्म है):
\[
1 + 3\varphi + 3\bar\varphi - 4 + 0 = 1 + 3 - 4 = 0,
\]
जिसमें $\varphi + \bar\varphi = 1$ का उपयोग हुआ है। सारणी हर लेखा-परीक्षा में उत्तीर्ण
होती है: वह $A_5$ की \omterm{def:b3:representations:table}{वर्ण सारणी} है।
\end{solution}
